%% kapa: exam=1 round=25 year=2014 period=1 form=A subject=civil_law subject_ko=민법 seq=35 number=35
%% source: Q-Net 감정평가사 기출문제 (cst00309) artlSeq=5101207 file=extracted/1008/5101207/2014년 제25회 감정평가사 1차시험 기출문제/1교시(A형).hwp
%% conversion: kordoc 4.24.51 -> tools/gen1.py (auto); review=proofread-sonnet; flags=none
\begin{question}{35}
전세권에 관한 설명으로 옳은 것은? (다툼이 있으면 판례에 의함)
\begin{choices}
  \item 전세권자는 목적부동산의 사용ㆍ수익의 방법에 어떠한 제한도 받지 않는다.
  \item 전세권설정자는 인도한 목적부동산을 사용ㆍ수익에 적합한 상태로 유지할 적극적 의무가 있다.
  \item 전세권의 존속기간 중 전세금반환청구권을 전세권과 분리하여 확정적으로 양도할 수 있다.
  \item 甲소유에 속한 대지와 건물 중 그 건물에 乙이 전세권을 취득한 후 그 대지가 丙에게 양도된 경우에 乙은 법정지상권을 취득한다.
  \item 주로 채권담보목적으로 전세권이 설정되었더라도 장차 전세권자가 목적물을 사용ㆍ수익하는 것을 완전히 배제하지 않았다면, 그 전세권의 효력은 인정된다.
\end{choices}
\end{question}
\answer{5}
